\section{Conclusiones y limitaciones}

Un único bloque iterado ocho veces (0.97\,M de parámetros) iguala o supera a un
GPT de ocho capas (6.5\,M) en los niveles difíciles y, sobre todo, fuera de la
distribución, y con los mismos parámetros mejora con claridad al GPT de una capa
(P1, P3). Con dos y cuatro iteraciones queda al nivel del GPT de igual
profundidad o algo por debajo. Las iteraciones necesarias crecen con la
dificultad (P2) y, con menos o con más de las entrenadas, la exactitud cae: no
hay extrapolación en profundidad. Por dentro, el Looped forma la respuesta de
forma gradual, más tarde cuanto mayor es el nivel, con casi la misma atención en
cada iteración, mientras que el GPT la forma en la última capa (P4).

\paragraph{Limitaciones.}
Modelos pequeños y datos sintéticos, que pueden no trasladarse a modelos
grandes ni a texto natural. Los niveles 1 y 2
miden menos razonamiento nuevo: el 100\,\% de sus trazas de test ya aparece en
el entrenamiento con otro texto. Hiperparámetros iguales para todos los modelos, sin ajuste por profundidad.
